% tema: fisica

\documentclass[12pt,a4paper]{article}

% Pacotes
\usepackage[utf8]{inputenc}
\usepackage[T1]{fontenc}
\usepackage[brazil]{babel}
\usepackage{amsmath,amssymb}
\usepackage{geometry}
\usepackage{booktabs}
\usepackage{array}
\usepackage{enumitem}
\usepackage{fancyhdr}
\usepackage{xcolor}
\usepackage{graphicx}

% Configurações
\geometry{
    left=2.5cm,
    right=2.5cm,
    top=2.5cm,
    bottom=2.5cm
}

\setlength{\parskip}{6pt}
\setlength{\parindent}{0pt}

\pagestyle{fancy}
\fancyhf{}
\lhead{Fundamentos de Mecânica}
\rhead{Avaliação Diagnóstica -- 2026}
\cfoot{\thepage}

\title{
    \textbf{Resolução da Avaliação Diagnóstica}\\
    \large Fundamentos de Mecânica (4300151)
}
\author{Camila Leite}
\date{2026}

\begin{document}

\maketitle

\tableofcontents
\newpage

\section{Parte I -- Questões de múltipla escolha}

\subsection{Questão 1 -- Movimento uniforme}

Dados:
\[
x_0=20\,\mathrm{m},\qquad
x_f=0\,\mathrm{m},\qquad
\Delta t=10\,\mathrm{s}
\]

A velocidade é:
\[
v=\frac{\Delta x}{\Delta t}
=\frac{0-20}{10}=-2\,\mathrm{m/s}
\]

A equação horária:
\[
x(t)=x_0+vt
\]

\[
\boxed{x(t)=20-2t}
\]

Resposta: alternativa C.

\subsection{Questão 2 -- Velocidade}

\[
v_m=\frac{\Delta S}{\Delta t}
\]

Intervalo de 0 a 10 s:
\[
v_1=\frac{50}{10}=5\,\mathrm{m/s}
\]

Intervalo de 10 a 20 s:
\[
v_2=0
\]

Intervalo de 20 a 40 s:
\[
v_3=\frac{-50}{20}=-2,5\,\mathrm{m/s}
\]

O maior módulo ocorre entre 0 e 10 s.

Resposta: alternativa E.

\subsection{Questão 3 -- Área sob o gráfico}

\[
\Delta x=\int_0^{1,6}v(t)\,dt
\]

Aproximando a área por trapézios:
\[
A_i=\frac{v_i+v_{i+1}}{2}\Delta t
\]

Com $\Delta t=0,2\,\mathrm{s}$:
\[
\Delta x\approx0,7\,\mathrm{m}
\]

Resposta: alternativa A.

\subsection{Questão 4 -- Lançamento vertical}

\[
v_f^2=v_0^2-2gh
\]

\[
0=v_0^2-2(10)(0,2)
\]

\[
v_0=\sqrt{4}
\]

\[
\boxed{v_0=2\,\mathrm{m/s}}
\]

Resposta: alternativa A.

\subsection{Questão 5 -- Queda livre}

\[
v_f^2=v_0^2+2gh
\]

\[
v_f=\sqrt{2(10)(5)}
\]

\[
v_f=10\,\mathrm{m/s}
\]

Conversão:
\[
v_f=10\cdot3,6
\]

\[
\boxed{v_f=36\,\mathrm{km/h}}
\]

Resposta: alternativa A.

\subsection{Questão 6 -- Encontro de automóveis}

\[
x_I=80t
\]

\[
x_P=\frac{8t^2}{2}
\]

Igualando:
\[
80t=4t^2
\]

\[
t=20\,\mathrm{s}
\]

Resposta: alternativa B.

\subsection{Questão 7 -- Energia cinética}

\[
E_m=E_c+E_p
\]

\[
E_c=E_m-E_p
\]

Nos extremos:
\[
E_c=0
\]

Na posição de equilíbrio, a energia cinética é máxima.

Resposta: alternativa C.

\subsection{Questão 8 -- Frenagem}

Velocidade após o tempo adicional de reação:
\[
v=14+1(1)=15\,\mathrm{m/s}
\]

Distância adicional durante a reação:
\[
\Delta x=14(1)+\frac{1(1)^2}{2}
\]

\[
\Delta x=14,5\,\mathrm{m}
\]

Distância de frenagem do motorista atento:
\[
d_A=\frac{14^2}{2(5)}=19,6\,\mathrm{m}
\]

Distância de frenagem do motorista desatento:
\[
d_D=\frac{15^2}{2(5)}=22,5\,\mathrm{m}
\]

Diferença total:
\[
\Delta d=14,5+(22,5-19,6)
\]

\[
\boxed{\Delta d=17,4\,\mathrm{m}}
\]

Resposta: alternativa E.

\subsection{Questão 9 -- Geometria da viga}

\[
L^2=x^2+y^2
\]

A distância entre o ponto médio e O:
\[
d_{MO}=
\sqrt{\frac{x^2}{4}+\frac{y^2}{4}}
\]

\[
d_{MO}=\frac{1}{2}\sqrt{x^2+y^2}
\]

\[
\boxed{d_{MO}=\frac{L}{2}}
\]

Resposta: alternativa A.

\subsection{Questão 10 -- Lançamento vertical}

No ponto de altura máxima:
\[
v=0
\]

Entretanto:
\[
a=-g
\]

A aceleração permanece orientada para baixo.

Resposta: alternativa C.

\newpage
\section{Parte II -- Problemas}

\subsection{Problema 1 -- Gráficos do movimento}

A aceleração média é:
\[
a_m=\frac{\Delta v}{\Delta t}
\]

De 0 a 1 minuto:
\[
a_m=\frac{12,5-0}{60}
\]

\[
\boxed{a_m=0,2083\,\mathrm{m/s^2}}
\]

De 2 a 3 minutos:
\[
a_m=\frac{0-12,5}{60}
\]

\[
\boxed{a_m=-0,2083\,\mathrm{m/s^2}}
\]

A distância total é obtida pela soma das áreas sob o gráfico:
\[
\boxed{\Delta s=3,375\,\mathrm{km}}
\]

\subsection{Problema 2 -- Malabarismo}

Altura máxima:
\[
h=2\,\mathrm{m}
\]

Velocidade inicial:
\[
v_0=\sqrt{2gh}
\]

\[
v_0=\sqrt{2(10)(2)}
\]

\[
\boxed{v_0=6,32\,\mathrm{m/s}}
\]

Tempo de subida:
\[
t_s=\frac{v_0}{g}
\]

\[
\boxed{t_s=0,632\,\mathrm{s}}
\]

Tempo total de voo:
\[
t_v=2t_s
\]

\[
\boxed{t_v=1,264\,\mathrm{s}}
\]

Intervalo entre lançamentos:
\[
\boxed{\Delta t=0,632\,\mathrm{s}}
\]

\subsection{Problema 3 -- Caçador e macaco}

Movimento da bala:
\[
x_b=v_0\cos\theta\,t
\]

\[
y_b=v_0\sin\theta\,t-\frac{gt^2}{2}
\]

Movimento do macaco:
\[
x_m=d
\]

\[
y_m=h-\frac{gt^2}{2}
\]

Condição horizontal do encontro:
\[
v_0\cos\theta\,t=d
\]

\[
t=\frac{d}{v_0\cos\theta}
\]

Como:
\[
\tan\theta=\frac{h}{d}
\]

Temos:
\[
\cos\theta=\frac{d}{\sqrt{d^2+h^2}}
\]

Logo:
\[
\boxed{
t=\frac{\sqrt{d^2+h^2}}{v_0}
}
\]

Na direção vertical:
\[
y_b=d\tan\theta-\frac{gt^2}{2}
\]

\[
y_b=h-\frac{gt^2}{2}
\]

Portanto:
\[
\boxed{y_b=y_m}
\]

A bala atinge o macaco porque ambos apresentam
o mesmo deslocamento vertical devido à gravidade.

\end{document}
